% Capítulo de la tesis. Estructura según la clase 7 y el material de apoyo de EP7000: responder las preguntas,
% verificar las hipótesis, contrastar con estudios previos, evaluar el método, implicancias, limitaciones y proyecciones.
\chapter{Discusión}
\label{sec:discusion}

Este capítulo interpreta los resultados del Capítulo~\ref{sec:resultados}. Primero responde las preguntas de investigación y verifica las hipótesis. Después contrasta los resultados con la literatura, evalúa el método y deriva las implicancias, las limitaciones y las proyecciones.

\section{Respuesta a las preguntas de investigación}
\label{sec:disc-preguntas}

Que la fragilidad bancaria importa para el spread soberano ya lo establecen \citet{Chari2024b}. Esta tesis partió de ese resultado y planteó dos preguntas.

\textbf{¿El riesgo económico doméstico afecta al spread soberano?} Se asocia con él, pero agrega poco a lo que ya explica la fragilidad. Es una asociación modesta ---una desviación estándar de riesgo de cola dentro del país, un spread un 10\% mayor, unos 17 puntos básicos--- entre los niveles de dos series que se mueven lento, y el diseño no permite leerla como un efecto causal. La lectura prudente es que la cola del crecimiento es un fundamento doméstico de largo plazo del spread en la mayoría de los países, no una fuente de información nueva sobre sus movimientos de corto plazo.

\textbf{Cuando ambos riesgos se agudizan en simultáneo, ¿el spread soberano se amplifica?} No hay evidencia de que lo haga. Los dos riesgos sí tienden a darse juntos, y ese es el resultado más firme del trabajo. Pero cuando coinciden, el spread no sube más que la suma de las dos asociaciones. Se descarta una amplificación de más de un 35\% de la suma, no una menor. La respuesta es una hipótesis nula que no se rechaza, no un resultado establecido: con trece países, el intervalo es ancho.

\section{Verificación de las hipótesis}
\label{sec:disc-hipotesis}

\textbf{H1 (canal bancario): se replica, con el alcance delimitado.} La asociación de \citet{Chari2024b} se replica con series construidas de forma independiente y con los controles de ese trabajo; con los seis controles domésticos de esta tesis desaparece en logaritmos. Esa caída admite dos lecturas. Es el patrón que se esperaría si esos controles fueran canales por los que opera la fragilidad ---el costo esperado del rescate deteriora las cuentas fiscales---, y en ese caso incluirlos subestima su coeficiente. Pero también calza con un fundamento macroeconómico común que mueve a la vez la fragilidad, las cuentas fiscales y el spread. La g-estimación apunta en la misma dirección con un método que no confunde un mediador con un control: la asociación de la fragilidad con el spread pasa por las cuentas fiscales, o comparte con ellas un fundamento común, y el panel no permite distinguir entre las dos lecturas.

El panel permite además decir qué no es esa asociación. No es una elasticidad común: las pendientes por país van de $-0{,}84$ a $+0{,}97$ y Turquía aporta el 70\% del coeficiente. Con trece países, basta excluir uno para cambiar la lectura de cada coeficiente. No es un resultado sobre toda la distribución de la fragilidad: la métrica tiene un piso en el que no informa. Con datos trimestrales no es precedencia temporal: si el spread se fecha con cuidado, la fragilidad y el spread se mueven dentro del mismo trimestre. Como $JLoss$ se construye con precios accionarios de los bancos, y esos precios y el spread soberano responden a las mismas noticias, ese movimiento conjunto no distingue una dirección. A frecuencia mensual, en cambio, la fragilidad se mueve primero: un mes antes que el spread, sin dirección inversa y sin depender de un país, aunque con una magnitud pequeña. Es lo que cabe esperar de una métrica hecha con precios accionarios de los bancos, que pueden recoger antes la misma información que el spread. Por eso no se lee como un efecto causal. Y los datos no informan si la asociación es de nivel o transitoria.

Los bancos individuales tampoco permiten ir más allá (Sección~\ref{sec:micro-prereg}). En un tercer preregistro el patrón va en la dirección del pasivo contingente ---la asociación se concentra en los bancos sistémicos---, pero ninguna de las seis hipótesis pasa la corrección de Holm, y la propiedad estatal o extranjera no ordena la transmisión como predice la lectura de rescate. Que la métrica sea robusta a la calibración del motor descarta que el resultado sea un artefacto de ella; que deje de ser significativa al exigir cuatro o cinco bancos cotizados por economía recuerda que, en algunas, $JLoss$ resume muy pocos bancos.

\textbf{H2 (canal de crecimiento): se cumple solo en niveles.} Sobre el modelo de \citet{Chari2024b}, el riesgo de cola del crecimiento agrega información dentro de la muestra, en el límite de lo que resiste la inferencia exigente. No aparece en los cambios a cuatro trimestres con el $GaR$ en pseudo--tiempo real (con el estimado en toda la muestra sí aparece con el \textit{wild bootstrap}, aunque no con Driscoll--Kraay), ni en las proyecciones locales, ni entre países, ni mejora el pronóstico fuera de la muestra. El coeficiente agrupado depende de Polonia, en parte por el efecto del logaritmo sobre spreads muy bajos, aunque el signo es positivo en diez de las trece economías. Y el \textit{wild bootstrap} no la vuelve evidencia limpia, porque no protege del sesgo de retroalimentación. La revisión econométrica matiza esa lectura en una dirección. Cuando cada economía responde distinto al ciclo financiero global, la asociación resiste la inferencia exigente y la exclusión de cada país, y el spread converge a una relación de largo plazo con el riesgo de cola. El horizonte refuerza esa lectura: con el riesgo de cola del año siguiente la asociación es mayor por punto porcentual ---un 4,7\%, en las trece economías, si bien por desviación estándar es parecida (un 8\%, unos 15 puntos básicos), porque ese riesgo varía menos---, aunque no sobrevive a la corrección sobre las 23 hipótesis primarias.

\textbf{H3 (amplificación): no se cumple.} La hipótesis con que se aprobó esta investigación queda sin respaldo. Ni el producto continuo, ni el contraste directo con indicadores de estrés conjunto, ni la regresión cuantílica en la cola alta del spread encuentran que la coincidencia de los dos riesgos vaya con un spread mayor que la suma de las dos asociaciones. Una prueba de equivalencia acota lo que queda abierto. El signo negativo de la interacción en logaritmos viene de Polonia y no se lee como atenuación; en puntos básicos, la estimación puntual es positiva. Que los efectos sean aditivos en logaritmos equivale a que sean proporcionales en puntos básicos. Esa complementariedad es álgebra de la forma funcional, coherente con una intensidad de incumplimiento que responde a los dos riesgos, y no es amplificación. La interacción positiva en cambios a cuatro trimestres y en puntos básicos, como la interacción en primeras diferencias de una versión anterior, refleja la forma proporcional que se ve en niveles, no una amplificación adicional. La coincidencia tampoco ensancha la distribución de los cambios del spread: en una prueba preregistrada a frecuencia mensual, cada riesgo por separado la ensancha, pero su coincidencia no la ensancha más (Sección~\ref{sec:quant-prereg}).

Lo que sí se sostiene es la simultaneidad, un hecho sobre los riesgos, no sobre el spread. Bancos frágiles y cola adversa del crecimiento coinciden más de lo que se esperaría por azar ---1,2 veces lo esperado descontado el calendario común, 1,9 veces frente a la independencia pura---, en las trece economías, y los episodios conjuntos duran más que por azar. Hay además un orden: la fragilidad bancaria antecede al deterioro de la cola del crecimiento, y no al revés, sin depender de un país ni de los precios accionarios que comparten las dos métricas. La lectura tiene dos límites. La simultaneidad viene de pocos episodios comunes a muchas economías: tres de cada cuatro observaciones conjuntas son de 2020 en adelante, aunque en 2020--2022 lo observado iguala lo esperado por el calendario y el exceso viene de 2015--2019 y de 2023--2026. Y la precedencia es pequeña: con tres o cuatro rezagos ya no resiste el \textit{wild bootstrap}, y no se puede fechar dentro del trimestre como se hizo con el spread, porque el $GaR$ es trimestral.

\textbf{H4 (deuda y riesgo de crecimiento): se cumple dentro de la muestra, sin confirmación fuera de ella.} La prueba preregistrada con datos trimestrales encontró un resultado firme y dos nulos. Con la inferencia conjunta sobre los cuatro preregistros, el resultado firme se sostiene para el riesgo de cola del año siguiente, no para el del trimestre siguiente (Sección~\ref{sec:inferencia-conjunta}). Fuera de la muestra no se confirma: en Hungría y Corea el signo se repite, pero la magnitud no coincide y la potencia es mínima, y un instrumento de términos de intercambio no muestra la intensificación (Sección~\ref{sec:n1n2}). Según las consecuencias fijadas de antemano, es una regularidad de la muestra original. Un modelo de límite fiscal calibrado tampoco lo explica: en logaritmos predice el signo contrario y, en puntos básicos, una interacción positiva mucho mayor que la que muestran los datos (Sección~\ref{sec:modelo-resultados}). Donde hay más deuda, una cola adversa del crecimiento pesa más en el spread: es el canal de la capacidad de pago, en el que el crecimiento futuro importa en proporción a la deuda que hay que servir. Las extensiones de H4 no se confirman: ni el \textit{doom loop} completo ni la concentración en la deuda en moneda extranjera.

\textbf{Otras condiciones: ninguna modifica la transmisión de forma demostrable.} La adecuación de reservas no atenúa la transmisión de la fragilidad. Y el efecto directo de la fragilidad no se separa del canal fiscal. También se examinaron la resiliencia de la política económica, el respaldo externo y quién tiene la deuda. Esta última se contrastó con una prueba preespecificada ---la participación de no residentes en la deuda del gobierno---, y según el criterio declarado no modera ninguno de los términos. La consecuencia, fijada de antemano, es que la interacción positiva que aparecía en una agrupación \textit{ex post} de once economías fuera de las crisis ($+0{,}034$; $+1{,}8$ puntos básicos) pasa a exploratoria. El dato descriptivo lo refuerza: esa agrupación no corresponde a quién tiene la deuda, porque Polonia es la cuarta economía del panel en participación extranjera. El respaldo externo medido por país y fecha no se asocia con una transmisión menor, y un estudio de eventos diario lo confirma con un diseño mejor identificado: las líneas \textit{swap} de la Reserva Federal de 2020 bajaron el spread al anunciarse, pero ninguna línea desacopló el spread de los bancos más que en las economías sin respaldo. Las asimetrías por régimen de crisis no resisten los errores agrupados por país, la permutación ni la corrección por comparaciones múltiples. El coeficiente de la fragilidad es menor donde hay más reservas, de forma estable al rezagarlas y más nítida en puntos básicos; pero en logaritmos depende de Polonia, no pasa la corrección por comparaciones múltiples y la cuenta corriente y el tipo de cambio real muestran interacciones de fuerza similar. La analogía del seguro ---un resguardo que abarata la prima mientras el riesgo es idiosincrásico y deja de hacerlo cuando el choque es común--- puede quedar como lectura del contraste por régimen del riesgo de cola, no como hallazgo. En todos los casos es ausencia de evidencia con intervalos anchos, no evidencia de que esas condiciones sean irrelevantes.

\section{Contraste con estudios previos}
\label{sec:disc-literatura}

\textbf{\citet{Chari2024b}.} El coeficiente de la fragilidad, entre $0{,}107$ y $0{,}164$ con los controles de ese trabajo y $0{,}206$ con el CDS, es del orden del que ellos reportan. El resultado concuerda, y ese acuerdo valida el panel, porque las series se construyeron de forma independiente. La diferencia está en el alcance: la asociación no sobrevive a los seis controles domésticos de esta tesis, depende de Turquía y no se distingue del sesgo de retroalimentación. Tampoco es la conclusión de esta tesis, sino su punto de partida.

\textbf{\citet{FarhiTirole2018b}.} El bucle se ve en los datos desde el otro extremo: el spread rezagado predice la fragilidad, porque un spread más alto deprecia los bonos soberanos que los bancos tienen en su balance y eleva su probabilidad de incumplimiento. La razón de fondo de la dificultad de identificar es la que motiva la hipótesis: el bucle hace que la relación opere en ambas direcciones y con mucha persistencia, y un coeficiente de nivel del tamaño observado es lo que produciría esa sola retroalimentación. En cambio, la amplificación que sugiere su multiplicador no aparece como $\beta_3>0$. Como anticipa el marco teórico, $\beta_3=0$ no contradice ese modelo: es su versión proporcional.

\textbf{\citet{BernankeGertlerGilchrist1996}.} La precedencia de la fragilidad sobre el riesgo de cola es el orden que predice el acelerador financiero, en el que el deterioro de los balances bancarios contrae el crédito y, con él, la actividad. Calza con que el triángulo del mecanismo conceptual tenga un lado más activo que el otro.

\textbf{\citet{ABG2019b} y \citet{AdrianEtAl2022}.} La mayor asociación con el riesgo de cola del año siguiente es el patrón que anticipa la estructura temporal del riesgo de cola: la vulnerabilidad que acumulan las condiciones financieras se manifiesta en el mediano plazo, y ese es el riesgo que más se asocia con el spread.

\textbf{Modelos de incumplimiento soberano.} El signo de H4 es el que sugieren \citet{Arellano2008} y \citet{AguiarGopinath2006}, pero el modelo de límite fiscal calibrado no explica el resultado en logaritmos, y la magnitud también falla: el spread responde al riesgo de cola diecisiete veces menos de lo que implicaría el modelo. Hay tres razones, no excluyentes, para esperarlo. El $GaR$ mide el riesgo de cola con error, lo que atenúa sus coeficientes. Una parte grande del spread es una prima global que no depende de la probabilidad de incumplimiento de cada economía. Y los inversionistas pueden no valorar la cola del crecimiento que estima el $GaR$ como lo hace el modelo. Según la consecuencia fijada en el preregistro, el modelo no explica la interacción en logaritmos. Las candidatas son extensiones que el modelo no incluye:
\begin{itemize}
    \item una recuperación que cae con la deuda \citep{Yue2010};
    \item inversionistas aversos al riesgo, cuya prima por la cola crece con la exposición \citep{Lizarazo2013};
    \item colas más gruesas: con la cola de Student, la interacción negativa se acerca a cero, sin volverse positiva.
\end{itemize}
Ninguna se contrastó aquí.

\section{Evaluación del método}
\label{sec:disc-metodo}

El método sirvió para acotar lo que se puede afirmar más que para identificar efectos. Varios resultados que parecían firmes en una primera estimación no sobrevivieron a las pruebas más exigentes: la interacción entre los dos riesgos, la asimetría por régimen de crisis, las reservas y, en parte, la intensificación del riesgo de cola por la deuda. La tesis los reporta así porque cada prueba se fijó antes de los datos y se publica salga como salga. Esa disciplina vale por sí misma. Con trece economías, reportar solo lo significativo produciría hallazgos que no se replicarían. Lo que queda en pie es menos, pero se sabe con qué grado de confianza queda.

El carácter estimado del $GaR$ no es el problema: al propagar su error, los errores estándar apenas crecen. La debilidad de la evidencia viene de la dinámica y del número de países, no de la medición del riesgo de cola.

\section{Implicancias}

Las implicancias de política son prudentes y siguen lo que el panel respalda.

\textbf{Vigilar los dos riesgos juntos se justifica.} No porque su coincidencia tenga un costo adicional demostrado, sino porque tienden a agravarse a la vez y porque la fragilidad bancaria antecede al deterioro de la cola del crecimiento. En ese sentido, una medida estructural de fragilidad bancaria informa más sobre el riesgo económico que viene que sobre el spread que viene.

\textbf{No hay base para una prima adicional por la simultaneidad.} La evidencia no respalda una exigencia regulatoria adicional para los episodios de doble vulnerabilidad por sobre la que corresponde a cada riesgo por separado. Tampoco permite descartar una moderada: el panel rechaza una amplificación de más de un 35\% de la suma, pero no una menor.

\textbf{El espacio fiscal podría importar para el riesgo económico.} Dentro de la muestra, y en dos preregistros, el mismo deterioro de la cola del crecimiento cuesta más en spread a una economía más endeudada. Si la regularidad se confirma con más economías, el riesgo de crecimiento debería vigilarse junto con la deuda. Por ahora no se confirma: la réplica en Hungría y Corea no reproduce su magnitud, el instrumento de términos de intercambio no la identifica y el modelo de límite fiscal predice en logaritmos el signo contrario (Secciones~\ref{sec:n1n2} y~\ref{sec:modelo-resultados}).

\textbf{No hay un umbral de vigilancia que los datos respalden.} La prueba preregistrada no encuentra un nivel de tensión financiera ni de deuda desde el cual la transmisión cambie, y los umbrales estimados país por país no son comunes ni transferibles entre economías (Sección~\ref{sec:umbrales-prereg}). Si la regularidad de la deuda se confirma, la regla que se desprende es continua, no de umbral: el mismo deterioro del riesgo de cola cuesta más en spread cuanto mayor es la deuda. Tampoco hay base para un semáforo de alerta temprana. Fuera de la muestra, ni la fragilidad, ni el riesgo de cola, ni la deuda anticipan los episodios de estrés soberano mejor que el azar, y la tensión financiera, la señal más informativa, no alcanza a ser útil. En los episodios de tensión, los riesgos domésticos pierden peso relativo en el spread: la vigilancia doméstica parece servir más en tiempos normales que en medio de una crisis.

\textbf{No hay base para recomendar reservas ni para diferenciar por quién tiene la deuda.} El término de las reservas es exploratorio, también en la prueba preregistrada, y no sostiene ninguna afirmación sobre lo que ocurriría si una economía las acumulara. La participación extranjera en la deuda no modera la transmisión según el criterio preespecificado, y el acceso a líneas de respaldo externo tampoco.

\textbf{Lo que el panel no permite.} No permite cuantificar cuánto bajaría el spread si cayera la fragilidad bancaria, porque las asociaciones no son efectos causales identificados. Tampoco permite usar ninguno de los dos riesgos para pronosticar el spread del trimestre siguiente.

\section{Limitaciones}

\begin{enumerate}
\item El corte transversal es de trece economías (catorce con Argentina, cuyo $GaR$ no es del todo homogéneo con el del resto). Con ese número de \textit{clusters}, en niveles los errores de Driscoll--Kraay rechazan una hipótesis nula cierta entre el 20\% y el 30\% de las veces; el \textit{wild cluster bootstrap} por país está bien calibrado, pero su potencia es de 35\% para la fragilidad y de 67\% para el riesgo de cola.
\item \textbf{Sesgo de retroalimentación.} El spread y la fragilidad son muy persistentes y se retroalimentan. En una simulación sin efecto directo, la regresión en niveles entrega en promedio $+0{,}066$ para la fragilidad (el dato es $0{,}099$) y $+0{,}022$ para el riesgo de cola (el dato es $0{,}032$); en diferencias, $+0{,}019$ para la fragilidad (el dato es $0{,}037$). Es un problema de identificación, no de errores estándar: ni el rezago, ni la especificación dinámica, ni las variables instrumentales disponibles, ni el \textit{wild bootstrap} lo resuelven.
\item \textbf{El piso de $JLoss$.} La métrica se amontona en $2{,}16$, el techo de la malla de pérdidas: el 23,7\% de las observaciones está a menos de un 2\% del piso, y más de la mitad en Malasia y Perú. En el piso la variable no informa, de modo que los coeficientes de la fragilidad describen los episodios fuera de él. $JLoss$ es además ordinal por la calibración del motor, y la asociación es sensible a una malla más ancha. Con la distancia al incumplimiento de Merton, que no satura, la asociación es más nítida ($p$ de \textit{wild bootstrap} de $0{,}055$ en logaritmos y $0{,}048$ en puntos básicos) y la interacción sigue siendo nula en logaritmos: la saturación atenúa la fragilidad, pero no explica la ausencia de amplificación.
\item \textbf{El spread es un promedio trimestral.} El EMBI del panel promedia los días del trimestre, y $JLoss$ usa precios de cierre. Con datos trimestrales se puede descartar que la fragilidad anteceda al spread de un trimestre al siguiente, pero no ordenar los movimientos dentro del trimestre; la serie mensual de la distancia al incumplimiento lo hace y encuentra que la fragilidad se mueve un mes antes (Sección~\ref{sec:mensual}), pero no existe un $JLoss$ mensual y el $GaR$ sigue siendo trimestral. Las 40 observaciones empalmadas del FMI no tienen dato diario. No se pudo verificar el rezago de publicación de los balances: ``mismo trimestre'' lo es para el precio y la volatilidad, no necesariamente para el balance.
\item \textbf{Pendientes heterogéneas.} La homogeneidad de pendientes se rechaza con claridad para la fragilidad y apenas para el riesgo de cola. Cada coeficiente agrupado depende de una economía: el de la fragilidad, de Turquía, y el del riesgo de cola, de Polonia.
\item La simultaneidad proviene de pocos episodios comunes, concentrados desde 2020, y su magnitud depende de la hipótesis nula: 1,2 veces lo esperado descontado el calendario común, 1,9 veces frente a la independencia pura. La precedencia de la fragilidad sobre el riesgo de cola resiste la exclusión de cada país y el $GaR$ sin bloque bursátil, pero con tres o cuatro rezagos ya no resiste el \textit{wild bootstrap} y no se puede fechar dentro del trimestre.
\item La prueba de la participación extranjera está preespecificada en el encabezado del código, no en un registro externo fechado, y el repositorio no lo acredita: el archivo se registró por primera vez el 21 de octubre de 2026, junto con sus resultados; tres cuartas partes de la varianza del moderador son entre países, de modo que compara en lo esencial trece economías. El moderador mide la tenencia extranjera de la deuda total del gobierno, y el EMBI es el spread de la deuda externa en dólares; la prueba sería más directa con la participación extranjera en la deuda en moneda local, que la base cubre de forma incompleta.
\item El acceso a respaldo externo no es exógeno: las líneas se piden y se otorgan según el riesgo, y solo seis economías lo tuvieron.
\item Los moderadores de la resiliencia no estaban preespecificados; las reservas, la deuda y el balance fiscal son series anuales interpoladas.
\item La forma funcional no está completamente determinada: los datos prefieren los logaritmos a los niveles, pero la transformación de Box--Cox óptima es intermedia.
\item La comparación fuera de la muestra no incluye efectos de tiempo y usa el $GaR$ sin reestimarlo recursivamente: es evidencia débil en cualquier dirección.
\item Los controles domésticos son interpolaciones de series anuales y plausiblemente median el mecanismo; faltan la calificación soberana histórica y el PIB per cápita que usan \citet{Chari2024b}.
\item La cobertura del EMBI es desigual: Filipinas aporta 12 trimestres, e Indonesia y Sudáfrica 32 cada una, 20 de ellos de la serie del FMI.
\end{enumerate}

\section{Proyecciones}
\label{sec:disc-proyecciones}

Las limitaciones anteriores marcan las oportunidades de investigación: más economías, variación exógena en la deuda y modelos con aversión al riesgo. El Capítulo~\ref{sec:agenda-futura} las desarrolla.
